为了证明 $\Th{Q}$ 能表示所有可计算函数，
我们需要一个精确的可计算性模型，作为证明的基础。
当然，可计算性模型有许多种，例如图灵机。
有一个在历史上扮演了重要角色的模型——它是最早被提出的模型之一，
也是 G\"odel 本人使用的模型——那就是\textbf{递归函数}。
递归函数是一类算术函数——即定义域和值域都是自然数——
它们可以从几个基本函数出发，通过若干种运算来定义。
基本函数是 $\Zero$、$\Succ$ 和投影函数。
运算是\textbf{复合}、\textbf{原始递归}和\textbf{无界搜索}。
复合就是将函数“链接”在一起的一般形式：先应用一个函数，
再将另一个函数应用于其结果。
原始递归从两个已定义的函数 $g$、$h$ 定义一个新函数~$f$，
规定 $f$ 在 $0$ 处的值由~$g$ 给出，
而在任意数 $n+1$ 处的值由 $h$ 作用于 $f(n)$ 给出。
仅用这两条原则就能定义的函数称为\textbf{原始递归的}。
一个关系是原始递归的，当且仅当其特征函数是原始递归的。

结果表明，一大批有趣的函数和关系都是原始递归的（例如加法、乘法、指数、整除性），
而且我们可以利用有界量化和有界极小化等原则，从已有的原始递归函数和关系出发，
定义出新的原始递归函数和关系。
特别地，这使我们能够说明，可以用原始递归的方式处理数的\emph{序列}。
也就是说，存在一种将数的序列“编码”为单个数的方法，使得我们能够计算第 $i$ 个元素、
序列长度、两个序列的连接等等，
而所有这些计算都由作用于这些编码的原始递归函数来完成。

为了得到部分递归函数，我们最终在复合和原始递归之外，
又添加了基于\textbf{无界搜索}的定义方式。
为了得到全可计算函数，我们将无界搜索限制在\textbf{正则}函数上。
一个函数~$g(x, y)$ 是正则的，当且仅当对每个~$y$，
它至少在某个~$x$ 处取值为~$0$。
如果 $g$ 是正则的，那么使得 $g(x, y) = 0$ 的最小 $x$ 总是存在的，
并且只需依次计算 $g(0, y)$、$g(1, y)$ 等等的值，
直到其中有一个等于 $0$，即可找到。
由此得到的函数~$f(y) = \umin{x}{g(x, y) = 0}$ 就是由~$g$ 经无界搜索定义的函数。
它始终是全函数且可计算的。
如此得到的函数集合称为\textbf{一般递归函数}。
Church-Turing（丘奇-图灵）论题的一个版本断言，可计算的算术函数恰好就是一般递归函数。